\section{Introduction}
3D scene understanding supports question answering, object description, and language grounding in spatial observations. Models such as LLaVA-3D, 3D-LLaVA, and LEO~\citep{zhu2024llava3d,deng20253dllava,huang2024leo} connect multi-view or point-cloud representations to a language model (LM). These 3D large multimodal models (3D LMMs) process hundreds to thousands of spatial tokens, increasing prefill computation and decoding key--value (KV) storage. We study token reduction for efficient inference without retraining these models.

Existing methods address two coupled choices: which tokens to retain and when to remove the rest. FastV uses last-prompt attention to rank visual tokens~\citep{chen2024fastv}. Depth-aware approaches also organize removal across layers: PyramidDrop exploits depth-dependent redundancy, while HiMAP uses information flow to guide hierarchical pruning~\citep{yuan2024pyramiddrop,yin2025himap}. Alongside these scheduling strategies, 3DZip uses spatial structure and feature diversity to compress 3D inputs~\citep{baek2026zip}. Together, these approaches motivate studying token selection and pruning depth jointly.

Suppose deleting the top-ranked spatial tokens hurts answer quality more at one layer than at another. One explanation is that the model needs those tokens most at the first layer. Another is that the first layer simply ranks a different, more useful set of tokens highly. Applying the same scoring rule at both layers does not distinguish these explanations. We therefore separate the layer that supplies the ranking from the layer where we delete the tokens. We then ask a different question about inference: at a given computational budget, which schedule works better when we \emph{keep} the top-ranked set?

In two Llama-family 7B models, one using image-derived tokens and the other superpoints, attention from the last prompt token to spatial tokens rises sharply from L13 to L14. This observation motivates \textbf{GeoGate}: run the full sequence through zero-indexed block L14, use FastV's ranking from that forward pass, and keep the selected spatial tokens before L15. GeoGate preserves all text tokens and original rotary positions and compacts every layer's prompt KV cache before decoding. ScanQA validation informed the choice of L14 and the default 10\% retention budget.

Our contributions are:
\begin{enumerate}
\item \textbf{A fixed-set test of deletion depth.} For tokens selected at L14 or L20, deleting the same set hurts exact match (EM) most at L8, then at L14, and least at L20. The L8-selected set shows no clear ordering. The peak observed when each layer selects its own tokens does not identify a peak due to deletion depth alone.
\item \textbf{A cost-matched test of late retention.} On LLaVA-3D ScanQA validation, keeping 128 spatial tokens at L14 yields 1.22 and 1.03 more EM points than selected L2 and L8 schedules at matched analytical LM cost. The layer sweep does not distinguish L14 from neighboring late settings.
\item \textbf{Limits across tasks and cost measures.} At 10\% retention, 3D-LLaVA loses 0.95 EM points on SQA3D validation. For LLaVA-3D, 42.1\% fewer analytical LM FLOPs translate to a 0.8--1.4\% reduction in measured application time on one H100.
\end{enumerate}
